% =========================================================
% CONSTRUCCIÓN CON EL TABLERO KANBAN (tarea O3-02)
% QUÉ VA: evidencia de que el plan de desarrollo (Tabla tab:ob2_tarjetas) se
%         ejecutó con Kanban: tablero real, fechas de cada tarjeta y desviaciones.
% DATOS:
%   - Captura del tablero real (GitHub Projects o Trello) en una fecha de corte.
%   - Fecha de inicio y fin de cada tarjeta: sale de git log de los dos
%     repositorios (primer y último commit asociado a la tarjeta).
%   - Desviaciones: tarjetas divididas (p. ej., odontograma en varias entregas),
%     tarjetas nuevas surgidas al integrar frontend y backend, cambios de orden.
% NO VA: pruebas, nombres de paquetes, detalles de código.
% EXTENSIÓN: 1 figura + 1 tabla + 1 párrafo de desviaciones.
% =========================================================

\subsubsection{Construcción}

La construcción siguió las quince tarjetas del plan de desarrollo (Tabla~\ref{tab:ob2_tarjetas}) en el tablero Kanban, con un límite de una tarjeta en desarrollo. Cada tarjeta abarcó el servidor y la interfaz del módulo. La Figura~\ref{fig:ob3_tablero} muestra el tablero al 22 de octubre de 2026, y la Tabla~\ref{tab:ob3_tarjetas}, las fechas de cada tarjeta.

\begin{figure}[H]
\centering
\IfFileExists{img/Ob3_Tablero_Kanban.png}{\includegraphics[width=\linewidth]{img/Ob3_Tablero_Kanban.png}}{\fbox{\parbox{0.9\linewidth}{\centering Pendiente: captura del tablero real en img/Ob3\_Tablero\_Kanban.png}}}
\caption{Tablero Kanban de la construcción}
\label{fig:ob3_tablero}
\end{figure}

% Fechas: primer y último commit de la tarjeta en los repositorios del
% backend, del frontend y del agente (git log). Días = días calendario.
\begingroup
\fontsize{9}{11}\selectfont
\renewcommand{\arraystretch}{1.2}
\setlength{\LTpre}{0pt}\setlength{\LTpost}{6pt}
\setlength{\tabcolsep}{4pt}
\begin{xltabular}{\linewidth}{|>{\raggedright\arraybackslash}X|c|c|c|}
\caption{Ejecución de las tarjetas del plan de desarrollo}
\label{tab:ob3_tarjetas} \\
\hline
\textbf{Tarjeta} & \textbf{Inicio} & \textbf{Fin} & \textbf{Días} \\ \hline
\endfirsthead
\hline
\textbf{Tarjeta} & \textbf{Inicio} & \textbf{Fin} & \textbf{Días} \\ \hline
\endhead
1. Identidad y autenticación & 08/10 & 11/10 & 4 \\ \hline
2. Profesionales y pacientes & 09/10 & 12/10 & 4 \\ \hline
3. Historia clínica y antecedentes & 09/10 & 12/10 & 4 \\ \hline
4. Agenda y horarios & 09/10 & 19/10 & 11 \\ \hline
5. Atención clínica y CIE-10 & 09/10 & 20/10 & 12 \\ \hline
6. Odontograma & 10/10 & 12/10 & 3 \\ \hline
7. Tratamientos & 17/10 & 19/10 & 3 \\ \hline
8. Recetas & 17/10 & 19/10 & 3 \\ \hline
9. Cobros y comprobante & 21/10 & 22/10 & 2 \\ \hline
10. Auditoría & 21/10 & 22/10 & 2 \\ \hline
\textbf{11. Agente: núcleo, aviso, menú y derivación} & 18/10 & 19/10 & 2 \\ \hline
\textbf{12. Agente: agendar y reprogramar} & 18/10 & 19/10 & 2 \\ \hline
\textbf{13. Agente: interpretación del texto libre y urgencias} & 19/10 & 19/10 & 1 \\ \hline
\textbf{14. Agente: recordatorios} & 18/10 & 19/10 & 2 \\ \hline
15. Respaldos y exportación & --- & --- & --- \\ \hline
\multicolumn{4}{|l|}{\cellcolor{gray!10}\textbf{Tarjetas añadidas durante la construcción}} \\ \hline
Formulario 033: vista e impresión & 19/10 & 20/10 & 2 \\ \hline
Despliegue en el servidor y CI/CD & 20/10 & 21/10 & 2 \\ \hline
Gestión de usuarios, catálogos y estadísticas del director & 21/10 & 22/10 & 2 \\ \hline
\end{xltabular}
\endgroup

% PENDIENTE: completar la tarjeta 15 y la ficha clínica del paciente
% (historial inmediato, RF06) cuando se terminen.
Las tarjetas de la agenda y de la atención se extendieron en varias entregas: la agenda, por los horarios, la disponibilidad en bloques de 15 min y la duración variable de la cita; la atención, por las casillas del Formulario 033. Durante la construcción se añadieron tres tarjetas: la vista e impresión del Formulario 033, el despliegue con integración y entrega continuas, y las funciones del director (usuarios, catálogos administrables y estadísticas). Además, el cobro se habilitó también para el odontólogo desde la atención. Al 22 de octubre quedan pendientes los respaldos automáticos y la conexión del agente con WhatsApp.
